\documentclass[11pt]{article}
\usepackage[margin=1in]{geometry}
\usepackage{amsmath,amssymb,amsthm}
\usepackage[vlined,linesnumbered]{algorithm2e}
\LinesNotNumbered
\usepackage{enumitem}

\setlength{\parindent}{0pt}
\setlength{\parskip}{0.5em}

\begin{document}


\section*{1. Greedy Knapsack}

\subsection*{(a)}

\textbf{Algorithm:} Sort items by weight ascending. Pick the lightest items one by one until the next item would exceed capacity.

\begin{algorithm}[H]
\DontPrintSemicolon
\SetKwBlock{Begin}{begin}{}
\Begin{
\textsc{UnitValueKnapsack}$(W[1..n],\, C)$\;
\textbf{Sort} $W$ in nondecreasing order\;
$k \longleftarrow 0$;\quad $w \longleftarrow 0$\;
\While{$k < n$ \textbf{and} $w + W[k+1] \le C$}{
    $k \longleftarrow k + 1$\;
    $w \longleftarrow w + W[k]$\;
}
\Return{$k$}\;
}
\end{algorithm}

\textbf{Proof (by contradiction):}
Let greedy select items $g_1, \ldots, g_k$ with $W[g_1] \le \cdots \le W[g_k]$ (the $k$ lightest).
Let OPT select $m$ items; sort them by weight as $W[o_1] \le \cdots \le W[o_m]$.

Since greedy chose the $k$ lightest items overall, the $j$-th lightest item overall is at most the $j$-th lightest in any subset:
\[
\forall\, j \le m: \quad W[g_j] \le W[o_j]
\]
Suppose $m > k$. Then:
\[
\sum_{j=1}^{m} W[o_j] \ge \sum_{j=1}^{k+1} W[o_j] \ge \sum_{j=1}^{k+1} W[g_j] > C \quad \text{(greedy stopped at } k\text{)}
\]
This contradicts OPT being feasible. \quad $\boxed{k \ge m \implies \text{greedy is optimal}}$

\textbf{Runtime:} Sort $O(n \log n)$ + scan $O(n)$ = $\boxed{O(n \log n)}$

\subsection*{(b)}

\textbf{Counterexample:} $n = 3$, $C = 2$, $W = [1, 1, 2]$, $V = [1, 1, 3]$.

\[
\begin{array}{c|ccc}
i & 1 & 2 & 3 \\ \hline
W[i] & 1 & 1 & 2 \\
V[i] & 1 & 1 & 3
\end{array}
\]

\textbf{Greedy:} Sorted by weight: items $1, 2, 3$. Pick items $1, 2$ (weight $1 + 1 = 2 \le C$). Value $= 2$.

\textbf{OPT:} Item $3$ alone: weight $2 \le C$, value $= 3$.

\[
\boxed{\text{Greedy} = 2 < 3 = \text{OPT}}
\]

\subsection*{(c)}

\textbf{Counterexample:} $n = 2$, $C = 2$, $W = [1, 2]$, $V = [2, 3]$.

\textbf{Greedy:} Sorted by weight: items $1, 2$. Pick item $1$ (weight $1$). Item $2$ does not fit ($1 + 2 = 3 > 2$). Value $= 2$.

\textbf{OPT:} Item $2$ alone: weight $2 \le C$, value $= 3$.

\[
\boxed{\text{Greedy} = 2 < 3 = \text{OPT}}
\]


\section*{2. Modified Interval Scheduling}

\subsection*{(a)}

\textbf{Counterexample:} $I_1 = [0, 2)$ with $V_1 = 1$, \; $I_2 = [1, 3)$ with $V_2 = 2$, \; $I_3 = [3, 5)$ with $V_3 = 1$.

\textbf{EFT greedy:} Sort by finish time: $I_1, I_2, I_3$. Pick $I_1$ (finish $2$). $I_2$ conflicts ($S_2 = 1 < 2 = F_1$), skip. Pick $I_3$ ($S_3 = 3 \ge 2$). Value $= 1 + 1 = 2$.

\textbf{OPT:} $\{I_2, I_3\}$: compatible ($F_2 = 3 \le 3 = S_3$). Value $= 2 + 1 = 3$.

\[
\boxed{\text{Greedy} = 2 < 3 = \text{OPT}}
\]

\subsection*{(b)}

\textbf{Greedy fails:} $I_1 = [0, 1)$ with $V_1 = 1$, \; $I_2 = [0.5, 1.5)$ with $V_2 = 100$.

EFT picks $I_1$ (finish $1$), then $I_2$ conflicts ($S_2 = 0.5 < 1$). Value $= 1$. OPT: $\{I_2\}$, value $= 100$.

\textbf{Adapted algorithm:}
Since all intervals have unit length, $F[i] = S[i] + 1$, so sorting by start time is the same as sorting by finish time. We adapt the standard weighted interval scheduling DP:

\begin{algorithm}[H]
\DontPrintSemicolon
\SetKwBlock{Begin}{begin}{}
\Begin{
\textsc{UnitLengthScheduling}$(S[1..n],\, V[1..n])$\;
\textbf{Sort} intervals by $S[i]$\;
$DP[0] \longleftarrow 0$\;
\For{$i \longleftarrow 1$ \textbf{to} $n$}{
    $p[i] \longleftarrow$ largest $j \ge 0$ with $S[j] + 1 \le S[i]$ \tcp*{binary search}
    $DP[i] \longleftarrow \max\!\big(DP[i{-}1],\; V[i] + DP[p[i]]\big)$\;
}
\Return{$DP[n]$}\;
}
\end{algorithm}

\textbf{Correctness:} $p[i]$ finds the latest non-conflicting interval ($F[j] = S[j]+1 \le S[i]$).
$DP[i]$ = max value using intervals $\{1, \ldots, i\}$: either skip $i$ ($DP[i{-}1]$) or take $i$ ($V[i] + DP[p[i]]$).

\textbf{Runtime:} Sort $O(n \log n)$ + $n$ binary searches $O(\log n)$ + DP fill $O(n)$ = $\boxed{O(n \log n)}$

\newpage
\section*{3. Interval Coloring}

\textbf{Data structure:} Sorted event array + stack of available room numbers.

\begin{algorithm}[H]
\DontPrintSemicolon
\SetKwBlock{Begin}{begin}{}
\Begin{
\textsc{IntervalColoring}$(S[1..n],\, F[1..n])$\;
Create $2n$ events: $(S[i],\, \text{start},\, i)$ and $(F[i],\, \text{end},\, i)$ for each $i$\;
\textbf{Sort} events by time (end before start at equal times)\;
$R \longleftarrow 0$;\quad $\mathrm{free} \longleftarrow$ empty stack\;
\For{\textbf{each} event in sorted order}{
    \If{event is \textnormal{end} for interval $j$}{
        $\mathrm{free}.\textsc{Push}(\mathrm{color}[j])$\;
    }
    \Else(\tcp*[f]{start for interval $j$}){
        \If{$\mathrm{free} \ne \emptyset$}{
            $\mathrm{color}[j] \longleftarrow \mathrm{free}.\textsc{Pop}()$ \tcp*{reuse room}
        }
        \Else{
            $R \longleftarrow R + 1$\;
            $\mathrm{color}[j] \longleftarrow R$ \tcp*{new room}
        }
    }
}
\Return{$R$}\;
}
\end{algorithm}

\textbf{Correctness:}
Whenever a new room $R$ is opened at time $S[i]$, the stack is empty, so all $R-1$ existing rooms are in use. Each of those rooms holds an interval that started before $S[i]$ and has not yet ended, so it overlaps $S[i]$. Together with $I_i$ itself, at least $R$ intervals are active at time $S[i]$:
\[
\text{depth}(S[i]) \ge R
\]
So $R \le D = \max_t \text{depth}(t)$: we never open more than $D$ rooms. Conversely, any valid coloring needs $\ge D$ rooms, so $R \ge D$. Hence $R = D$. \qed

\textbf{Runtime:} Create events $O(n)$ + sort $O(n \log n)$ + sweep $O(n)$ stack ops = $\boxed{O(n \log n)}$

%\newpage
\section*{4. Interval Coloring with Small and Large Rooms}

\textbf{Lower bound:}
\begin{align*}
&D = \max_t d(t) \quad \text{(max depth of all intervals)} \\
&D_1 = \max_t d_1(t) \quad \text{(max depth of } B[i]=1 \text{ intervals)} \\
&R_S + R_L \ge D \quad \text{(total rooms } \ge \text{ max depth)} \\
&R_L \ge D_1 \quad \text{(large rooms } \ge \text{ max depth of large-only)} \\
&R_S + 2R_L = (R_S + R_L) + R_L \ge D + D_1
\end{align*}

\textbf{Algorithm:} Two-pass event sweep. Pass~1 computes $D$. Pass~2 assigns rooms using two stacks
$\mathrm{free}_L$, $\mathrm{free}_S$ and array $\mathrm{tag}[\cdot]$ tracking room type.

\begin{algorithm}[H]
\DontPrintSemicolon
\SetKwBlock{Begin}{begin}{}
\Begin{
\textsc{SmallLargeColoring}$(S[1..n],\, F[1..n],\, B[1..n])$\;
Create $2n$ events: $(S[i],\, \text{start},\, i)$ and $(F[i],\, \text{end},\, i)$ for each $i$\;
\textbf{Sort} events by time (end before start at equal times)\;
\BlankLine
\tcp{Pass 1: compute $D$}
$D \longleftarrow 0$;\quad $\mathrm{active} \longleftarrow 0$\;
\For{\textbf{each} event in sorted order}{
    \lIf{end}{$\mathrm{active} \longleftarrow \mathrm{active} - 1$}
    \lElse{$\mathrm{active} \longleftarrow \mathrm{active} + 1$;\quad $D \longleftarrow \max(D,\, \mathrm{active})$}
}
\BlankLine
\tcp{Pass 2: assign rooms}
$R_S \longleftarrow 0$;\quad $R_L \longleftarrow 0$\;
$\mathrm{free}_S, \mathrm{free}_L \longleftarrow$ empty stacks\;
\For{\textbf{each} event in sorted order}{
    \If{event is \textnormal{end} for interval $j$}{
        Push $\mathrm{color}[j]$ onto $\mathrm{free}_L$ or $\mathrm{free}_S$ based on $\mathrm{tag}[\mathrm{color}[j]]$\;
    }
    \ElseIf{$B[j] = 1$ \tcp*[f]{start, needs large room}}{
        \If{$\mathrm{free}_L \ne \emptyset$}{
            $\mathrm{color}[j] \longleftarrow \mathrm{free}_L.\textsc{Pop}()$ \tcp*{reuse large}
        }
        \ElseIf{$\mathrm{free}_S \ne \emptyset$}{
            $r \longleftarrow \mathrm{free}_S.\textsc{Pop}()$\;
            $\mathrm{tag}[r] \longleftarrow L$ \tcp*{promote small $\to$ large}
            $R_S \longleftarrow R_S - 1$;\quad $R_L \longleftarrow R_L + 1$\;
            $\mathrm{color}[j] \longleftarrow r$\;
        }
        \Else{
            $R_L \longleftarrow R_L + 1$\;
            $\mathrm{color}[j] \longleftarrow R_S + R_L$;\quad $\mathrm{tag}[\mathrm{color}[j]] \longleftarrow L$ \tcp*{new large}
        }
    }
    \Else(\tcp*[f]{start, $B[j]=0$, needs any room}){
        \If{$\mathrm{free}_S \ne \emptyset$}{
            $\mathrm{color}[j] \longleftarrow \mathrm{free}_S.\textsc{Pop}()$ \tcp*{reuse small}
        }
        \ElseIf{$R_S + R_L < D$}{
            $R_S \longleftarrow R_S + 1$\;
            $\mathrm{color}[j] \longleftarrow R_S + R_L$;\quad $\mathrm{tag}[\mathrm{color}[j]] \longleftarrow S$ \tcp*{new small}
        }
        \Else{
            $\mathrm{color}[j] \longleftarrow \mathrm{free}_L.\textsc{Pop}()$ \tcp*{reuse large (forced)}
        }
    }
}
\Return{$R_S + 2 R_L$}\;
}
\end{algorithm}

The idea: when a $B[j]=0$ interval needs a room and no small room is free, we open a new small room instead of handing it a large room. This saves large rooms for future $B[i]=1$ intervals. We only fall back to reusing a large room when we have already opened $D$ rooms and cannot open more.

\textbf{Proof:}

\textbf{(A) Total rooms $= D$:}
A new room is opened in two cases: (i) both stacks are empty, meaning every existing room is in use, so the current depth $\ge$ total rooms and total stays $\le D$; or (ii) $B[j]=0$ with $\mathrm{free}_S = \emptyset$ and total $< D$, so total stays $\le D$.
Since any valid coloring needs $\ge D$ rooms, we get $R_S + R_L = D$.

\textbf{(B) $R_L = D_1$:}
Since cost $= D + R_L$, minimizing cost is equivalent to minimizing $R_L$.

$R_L \ge D_1$: at a time when $D_1$ intervals with $B[i]=1$ overlap, each sits in a distinct large room.

$R_L \le D_1$: call a room \emph{tainted} the moment it first becomes large (created large or promoted). Only a $B[j]=1$ interval can taint a room (when $\mathrm{free}_L = \emptyset$). A $B[j]=0$ interval never taints: it either reuses a free small room, opens a new small room (when total $< D$), or reuses a room that is already large.

Consider the $r$-th tainting, caused by some $B[i]=1$ interval $I_i$. Since $\mathrm{free}_L = \emptyset$, all $r-1$ existing large rooms are occupied by intervals that overlap $S[i]$.
We claim each of those occupants has $B=1$. Indeed, the only way a $B[j]=0$ interval sits in a large room is through the forced-reuse case (total $= D$ and $\mathrm{free}_S = \emptyset$). But that means \emph{all} $D$ rooms were occupied, so no room was free at all, contradicting $\mathrm{free}_L \ne \emptyset$.

So at time $S[i]$ there are $r-1$ overlapping $B=1$ intervals plus $I_i$ itself, giving $d_1(S[i]) \ge r$ and hence $R_L \le D_1$.

$\therefore\; R_L = D_1$, \quad $\text{cost} = D + D_1$, matching the lower bound. \qed

\textbf{Runtime:} Two passes over $2n$ sorted events, each $O(n)$ + sort $O(n \log n)$ = $\boxed{O(n \log n)}$


\end{document}
